%% 1-5 ひき算は「のこり」「ちがい」「かわった 量」（10ページ、p.56〜65）
%% 第1巻のきまり：15 pt・行送り28 pt、分かち書き（文節の間に半角の空白）、漢字はすべてルビ、式は全角の－＝で書く

%% ---- この章だけで使う絵の部品（頭文字 \bkqe…） ----
\makeatletter
% たての ものさし  \bkqevruler{(x,y)}{目もり1つの長さ}{いちばん上の目もりの数}   (x,y) は 0の目もり。数字は5ごと  % noruby
\newcommand\bkqevruler[3]{\begin{scope}[shift={#1}]
  \fill[bkYellowBg, draw=bkYellow!80!black, line width=0.8pt, rounded corners=1mm] (-5,-2.5) rectangle (5,#3*#2+2.5);
  \foreach \k in {0,...,#3} {\pgfmathtruncatemacro\bkqe@f{mod(\k,5)==0}
    \ifnum\bkqe@f=1 \draw[bkInk, line width=0.8pt] (-5,\k*#2) -- (-0.8,\k*#2);
      \node[font=\sffamily\bfseries\footnotesize, text=bkInk, inner sep=0pt] at (2.4,\k*#2) {\k};
    \else\draw[bkInk, line width=0.5pt] (-5,\k*#2) -- (-2.6,\k*#2);\fi}
\end{scope}}
% いす（前から見た形）  \bkqechair{(x,y)}{高さ}{色}   (x,y) は 足もとの まん中  % noruby
\newcommand\bkqechair[3]{\begin{scope}[shift={#1}, scale=#2/11]\colorlet{bkqe@c}{#3}
  \fill[bkqe@c!75!black] (-3,0) rectangle (-2.2,4.8) (2.2,0) rectangle (3,4.8);
  \fill[bkqe@c!75!black] (-3.2,5) rectangle (-2.3,10.6) (2.3,5) rectangle (3.2,10.6);
  \fill[bkqe@c, rounded corners=0.8] (-3.4,7.6) rectangle (3.4,10.8);
  \fill[bkqe@c, rounded corners=0.6] (-3.8,4.4) rectangle (3.8,5.8);
\end{scope}}
% なくなった ものの あと（点線の まる）  \bkqegone{(x,y)}{半径}  % noruby
\newcommand\bkqegone[2]{\draw[bkGray, line width=0.8pt, dash pattern=on 1.6pt off 1.4pt] #1 circle (#2);}
% ブロック（すきまを あけて つむ 1こ）  \bkqeblock{(x,y)}{1辺}{色}  % noruby
\newcommand\bkqeblock[3]{\colorlet{bkqe@c}{#3}\filldraw[fill=bkqe@c, draw=bkqe@c!65!black, line width=0.6pt, rounded corners=0.5mm] #1 rectangle ++(#2,#2);}
% スピードメーター  \bkqemeter{(x,y)}{半径}{はり（0〜100）}  % noruby
\newcommand\bkqemeter[3]{\begin{scope}[shift={#1}]
  \filldraw[fill=bkInk!88, draw=bkGray, line width=1pt] (0,0) circle (#2);
  \foreach \k in {0,...,10} {\pgfmathsetmacro\bkqe@a{210-24*\k}
    \draw[white, line width=0.7pt] (\bkqe@a:#2*0.9) -- (\bkqe@a:#2*0.75);}
  \foreach \k/\v in {0/0,2/20,4/40,6/60,8/80,10/100} {\pgfmathsetmacro\bkqe@a{210-24*\k}
    \node[text=white, font=\sffamily\bfseries\fontsize{7.5}{8}\selectfont, inner sep=0pt] at (\bkqe@a:#2*0.55) {\v};}
  \pgfmathsetmacro\bkqe@n{210-2.4*#3}
  \draw[bkOrange, line width=1.4pt, line cap=round] (0,0) -- (\bkqe@n:#2*0.86);
  \fill[white] (0,0) circle (#2*0.1);
\end{scope}}
\makeatother

\bkchapter{1}{5}{ひき\R{算}{ざん}は「のこり」\\「ちがい」「かわった \R{量}{りょう}」}{\bklevel{1}}{ひき\R{算}{ざん}は「のこり」「ちがい」「かわった \R{量}{りょう}」}
\bkmichishirube{1-1,1-2,1-4}{1-6,1-7,1-18,2-2}

\begin{goal}
\begin{itemize}[label={\color{bkTheme}●}, leftmargin=1.4\zw, itemsep=0pt, topsep=0pt]
\item 「のこり」「ちがい」を、 ひき\R{算}{ざん}の しきに かける。
\item まえと あとで、 どれだけ かわったかを もとめられる。
\item 10の まとまりを くずして、 ひき\R{算}{ざん}が できる。
\end{itemize}
\end{goal}

\begin{hatena}
\begin{bkfigure}
\begin{tikzpicture}[x=1mm, y=1mm]
  \fill[white, rounded corners=3mm] (0,0) rectangle (122,53);
  \node[anchor=north west, font=\sffamily\bfseries\small, text=bkLeaf!80!black] at (4,51) {あさがおの たかさ};
  % きのうと きょうの あさがお（ものさしの 目もり 12 と 15）  % noruby
  \bkplant{(36,15)}{24}
  \bkplant{(66,15)}{30}
  \node[font=\sffamily\bfseries\small, text=bkInk] at (36,3) {きのう};
  \node[font=\sffamily\bfseries\small, text=bkInk] at (66,3) {きょう};
  \bkqevruler{(94,15)}{2}{16}
  \draw[bkBlue, line width=1pt, dash pattern=on 2.2pt off 1.6pt] (41.5,39) -- (88.6,39);
  \draw[bkAppleRed, line width=1pt, dash pattern=on 2.2pt off 1.6pt] (71.5,45) -- (88.6,45);
  \node[draw=bkBlue, line width=0.9pt, fill=white, rounded corners=1mm, font=\sffamily\bfseries\small, inner sep=1.5pt] at (24,39) {12};
  \node[draw=bkAppleRed, line width=0.9pt, fill=white, rounded corners=1mm, font=\sffamily\bfseries\small, inner sep=1.5pt] at (54,45) {15};
  \draw[bkTheme, line width=1pt, decorate, decoration={brace, amplitude=1.4mm, mirror}] (101,39) -- (101,45);
  \node[font=\sffamily\bfseries\Large, text=bkTheme] at (107.5,42) {？};
\end{tikzpicture}
\end{bkfigure}
\noindent きのうは \R{目}{め}もり 12、 きょうは \R{目}{め}もり 15。\\
きのうより、 どれだけ のびたかな？\\
のびた ぶんは、 どんな しきで もとめるのかな？
\end{hatena}

\clearpage
%% ---------------- p.57 やってみよう ----------------
\begin{yattemiyou}[おはじきと ブロックで しらべよう]
\begin{enumerate}[label={\bkcircled[bkGreen]{\arabic*}}, leftmargin=2.2\zw, itemsep=2pt, topsep=0pt]
\item おはじきを 5こ ならべて、 2こ とろう。 のこりは なんこ？\enspace\bkkakikomi\enspace こ
\end{enumerate}
\begin{bkfigure}
\begin{tikzpicture}[x=1mm, y=1mm]
  \foreach \x in {0,10,20} {\bkohajiki{(\x,0)}{3.2}{bkRoute2}}
  \foreach \x in {30,40} {\bkqegone{(\x,0)}{3.2}}
  \foreach \x in {62,72} {\bkohajiki{(\x,6)}{3.2}{bkRoute2}}
  \draw[bk arrow, bkGreen] (45,2) -- (56,5);
  \node[font=\sffamily\bfseries\small, text=bkGreen!70!black] at (67,-2) {2こ とる};
\end{tikzpicture}
\end{bkfigure}
\begin{enumerate}[label={\bkcircled[bkGreen]{\arabic*}}, leftmargin=2.2\zw, itemsep=2pt, topsep=0pt, start=2]
\item あかい おはじき 5こと、 あおい おはじき 3こを、 1つずつ ペアに しよう（1-1）。 ペアに ならない おはじきは なんこ？\enspace\bkkakikomi\enspace こ
\end{enumerate}
\begin{bkfigure}
\begin{tikzpicture}[x=1mm, y=1mm]
  \foreach \x in {0,10,20,30,40} {\bkohajiki{(\x,10)}{3.2}{bkAppleRed}}
  \foreach \x in {0,10,20} {\bkohajiki{(\x,0)}{3.2}{bkRoute2}}
  \foreach \x in {0,10,20} {\draw[bkGreen, line width=1.2pt] (\x,6.4) -- (\x,3.6);}
  \draw[bkOrange, line width=1pt, dash pattern=on 2pt off 1.5pt, rounded corners=2mm] (25.5,5) rectangle (44.5,15);
  \node[anchor=west, font=\sffamily\bfseries\small, text=bkOrange!80!black] at (47,10) {ペアに ならない？};
\end{tikzpicture}
\end{bkfigure}
\begin{enumerate}[label={\bkcircled[bkGreen]{\arabic*}}, leftmargin=2.2\zw, itemsep=2pt, topsep=0pt, start=3]
\item ブロックを 4こ つもう。 7こに なるまで つみたそう。 なんこ つみたしたかな？\enspace\bkkakikomi\enspace こ
\end{enumerate}
\begin{bkfigure}
\begin{tikzpicture}[x=1mm, y=1mm]
  \foreach \k in {0,...,3} {\bkqeblock{(0,\k*5.6)}{5}{bkBlue}}
  \foreach \k in {4,5,6} {\draw[bkGreen, line width=0.9pt, dash pattern=on 1.6pt off 1.4pt, rounded corners=0.5mm] (0,\k*5.6) rectangle ++(5,5);}
  \draw[bkBlue, line width=0.8pt] (7.5,0) -- (9,0) -- (9,22) -- (7.5,22);
  \node[anchor=west, font=\sffamily\bfseries\small, text=bkBlue!80!black] at (11,11) {はじめ 4こ};
  \draw[bkGreen!70!black, line width=0.8pt] (-2.5,0) -- (-4,0) -- (-4,38.8) -- (-2.5,38.8);
  \node[anchor=east, font=\sffamily\bfseries\small, text=bkGreen!70!black] at (-6,19.4) {7こに なるまで};
\end{tikzpicture}
\end{bkfigure}
\noindent ゆびで かぞえても いいよ。 なんども ためして みよう。
\end{yattemiyou}

\bkfukidashi{mina}{①は へった、 ②は くらべた、 ③は のびた。 ちがう ことなのに、 どれも おなじ けいさんで もとめられるよ！}

\clearpage
%% ---------------- p.58〜59 数学の道具 ----------------
\bksection{doguu}{\R{数学}{すう|がく}の \R{道具}{どう|ぐ}}

\bksubsection{1}{「のこり」の ひき\R{算}{ざん}}

\label{bk:1-5:nokori}クッキーが 5まい。 2まい たべると、 のこりは なんまい？

\begin{bkfigure}
\begin{tikzpicture}[x=1mm, y=1mm]
  \bkplate{(30,0)}{66}
  \foreach \x in {5,16,27} {\bkcookie{(\x,1.5)}{4.2}}
  \foreach \x in {42,53} {\bkqegone{(\x,1.5)}{4.2}}
  \draw[bkStem, line width=0.8pt] (1,-6.5) -- (31,-6.5);
  \draw[bkGray, line width=0.8pt] (38,-6.5) -- (57,-6.5);
  \node[anchor=north, font=\sffamily\bfseries\small, text=bkStem!90!black] at (16,-7.5) {のこり 3まい};
  \node[anchor=north, font=\sffamily\bfseries\small, text=bkGray] at (47.5,-7.5) {たべた 2まい};
\end{tikzpicture}
\end{bkfigure}
\noindent この ことを、 しきで こう かきます。
\bkbigshiki{5 － 2 ＝ 3}{5 ひく 2 は 3}
「－」は 「ひく」と よみます。 これを ひき\R{算}{ざん}と いいます。\bkidx{ひきざん}{ひき\R{算}{ざん}}\bkidx{ひく}{－（ひく）}
「あげる」「とる」で へった のこりも、 ひき\R{算}{ざん}で もとめます。\bkidx{のこり}{のこり}\par
しきは、 「はじめの かず － へった かず」の じゅんに かきます。 はじめに 5まい あって、 そこから 2まい へったので 5 － 2 です。

\bksubsection{2}{「ちがい」の ひき\R{算}{ざん}}

\label{bk:1-5:chigai}いすが 4つ、 こどもが 6\R{人}{にん}。 こどもは なん\R{人}{にん} おおい？

\begin{bkfigure}
\begin{tikzpicture}[x=1mm, y=1mm]
  \foreach \x/\c in {0/bkRoute2,12/bkAppleRed,24/bkGreen,36/bkRoute4,48/bkOrange,60/bkPurple} {\bkkid{(\x,12)}{13.5}{\c}{bkHair}}
  \foreach \x in {0,12,24,36} {\bkqechair{(\x,0)}{9}{bkStem!70}}
  \foreach \x in {0,12,24,36} {\draw[bkGreen, line width=1.2pt] (\x,11.4) -- (\x,9.6);}
  \draw[bkOrange, line width=1pt, dash pattern=on 2pt off 1.5pt, rounded corners=2mm] (42.5,10.5) rectangle (65.5,27.5);
  \node[anchor=west, font=\sffamily\bfseries\small, text=bkOrange!80!black] at (68,19) {ペアに ならない};
\end{tikzpicture}
\end{bkfigure}
\noindent ひとりずつ いすと ペアに すると（1-1）、 2\R{人}{にん} あまります。
\bkbigshiki{6 － 4 ＝ 2}{6 ひく 4 は 2}
くらべて おおい ぶんが 「ちがい」。 これも ひき\R{算}{ざん}です。\bkidx{ちがい}{ちがい}

%% （第11回：ここで ページを かえない）
\bksubsection{3}{「かわった \R{量}{りょう}」の ひき\R{算}{ざん}}

\label{bk:1-5:kawatta}\noindent\begin{minipage}[c]{0.52\linewidth}\centering
\begin{tikzpicture}[x=1mm, y=1mm]
  \bkplant{(8,8)}{24}
  \bkplant{(32,8)}{30}
  \node[font=\sffamily\bfseries\footnotesize, text=bkInk] at (8,-3.2) {きのう};
  \node[font=\sffamily\bfseries\footnotesize, text=bkInk] at (32,-3.2) {きょう};
  \bkqevruler{(56,8)}{2}{16}
  \draw[bkBlue, line width=1pt, dash pattern=on 2.2pt off 1.6pt] (13.5,32) -- (50.6,32);
  \draw[bkAppleRed, line width=1pt, dash pattern=on 2.2pt off 1.6pt] (37.5,38) -- (50.6,38);
  \fill[bkOrange!80] (51.2,32) rectangle (54,38);
  \draw[bkTheme, line width=1pt, decorate, decoration={brace, amplitude=1.3mm, mirror}] (63,32) -- (63,38);
  \node[anchor=west, font=\sffamily\bfseries\small, text=bkTheme] at (65,35) {3};
\end{tikzpicture}
\end{minipage}\hfill
\begin{minipage}[c]{0.45\linewidth}
12から 15まで、 \R{目}{め}もりを かぞえると 13、 14、 15の 3つ。 のびた ぶんは 3\R{目}{め}もりです。
\end{minipage}

\noindent ひき\R{算}{ざん}でも もとめられます。 \R{目}{め}もりを ブロックに して つんで みよう。

\noindent\begin{minipage}[c]{0.36\linewidth}\centering
\begin{tikzpicture}[x=1mm, y=1mm]
  \foreach \k in {0,...,11} {\bkqeblock{(0,\k*2.6)}{2.4}{bkBlue}}
  \foreach \k in {12,13,14} {\bkqeblock{(0,\k*2.6)}{2.4}{bkOrange}}
  \draw[bkBlue!80!black, line width=0.8pt, decorate, decoration={brace, amplitude=1.3mm, mirror}] (3.6,0) -- (3.6,31);
  \node[anchor=west, font=\sffamily\bfseries\footnotesize, text=bkBlue!80!black] at (6,15.5) {きのうの 12};
  \draw[bkOrange!80!black, line width=0.8pt, decorate, decoration={brace, amplitude=1.1mm, mirror}] (3.6,31.4) -- (3.6,38.8);
  \node[anchor=west, font=\sffamily\bfseries\footnotesize, text=bkOrange!80!black] at (6,35.1) {のびた ぶん};
  \node[font=\sffamily\bfseries\footnotesize] at (1.2,-3.5) {きょう 15};
\end{tikzpicture}
\end{minipage}\hfill
\begin{minipage}[c]{0.6\linewidth}
きょうの 15は、 きのうの 12と、 のびた ぶんを あわせた ものです。 だから、 15から きのうの 12を とると、 のこりが のびた ぶんに なります（「のこり」と おなじ かんがえ）。
\end{minipage}
\bkbigshiki{15 － 12 ＝ 3}{あと － まえ}
\bkhatenaref の こたえは、 3\R{目}{め}もり。 きのうより のびた ぶん、 まえより ふえた ぶんは、 「あと － まえ」の ひき\R{算}{ざん}。 これを 「かわった \R{量}{りょう}」と いいます。\bkidx{かわったりょう}{かわった \R{量}{りょう}}

\bksubsection{4}{10の まとまりを くずす ひき\R{算}{ざん}}

\label{bk:1-5:kuzusu}おはじきが 13こ。 5こ とると、 のこりは なんこ？

\noindent\begin{minipage}[c]{0.4\linewidth}
\begin{tikzpicture}[x=1mm, y=1mm]
  \bktenframe{(0,0)}{6.4}{10}{bkRoute2}{0}{bkRoute2}
  \foreach \i in {0,...,4} {\draw[bkRed, line width=1.2pt, line cap=round] (\i*6.4+1.6,8) -- (\i*6.4+4.8,11.2) (\i*6.4+1.6,11.2) -- (\i*6.4+4.8,8);}
  \foreach \x in {38,44,50} {\bkohajiki{(\x,3.2)}{2.5}{bkRoute2}}
  \node[anchor=north, font=\sffamily\bfseries\footnotesize] at (16,-1.5) {10の まとまり};
  \node[anchor=north, font=\sffamily\bfseries\footnotesize] at (44,-1.5) {ばら 3こ};
\end{tikzpicture}
\end{minipage}\hfill
\begin{minipage}[c]{0.58\linewidth}
\begin{enumerate}[label={\bkcircled[bkTheme]{\arabic*}}, leftmargin=2\zw, itemsep=0pt, topsep=0pt]
\item ばらの 3こでは たりない。
\item 10から 5を とって 5。
\item 5と 3で 8。
\end{enumerate}
\end{minipage}
\bkbigshiki{13 － 5 ＝ 8}{13 ひく 5 は 8}

\noindent 13は、 10の まとまりと ばら 3こ（1-2）。 ばらの 3こからは 5こ とれないので、 10の まとまりを くずして 5こ とります。 10 － 5 ＝ 5 です（5と 5で 10。 1-4の 10に なる くみ）。 のこった 5こと、 ばらの 3こを あわせて 8こです。

%% （第11回：ここで ページを かえない）
\noindent べつの やりかたも あります。 とる 5を、 3と 2に わけます（いくつと いくつ、 1-4）。

\begin{bkfigure}
\begin{tikzpicture}[x=1mm, y=1mm]
  % ① 13
  \bktenframe{(0,0)}{5}{10}{bkRoute2}{0}{bkRoute2}
  \foreach \x in {29,34,39} {\bkohajiki{(\x,2.5)}{1.9}{bkRoute2}}
  \foreach \x in {29,34,39} {\draw[bkRed, line width=1pt] (\x-1.6,0.9) -- (\x+1.6,4.1) (\x-1.6,4.1) -- (\x+1.6,0.9);}
  \node[font=\sffamily\bfseries\footnotesize, text=bkRed] at (38,9.5) {3こ とる};
  \draw[bk arrow, bkTheme] (45,5) -- (52,5);
  % ② 10
  \bktenframe{(55,0)}{5}{10}{bkRoute2}{0}{bkRoute2}
  \foreach \i in {3,4} {\draw[bkRed, line width=1pt] (55+\i*5+1,1) -- (55+\i*5+4,4) (55+\i*5+1,4) -- (55+\i*5+4,1);}
  \node[font=\sffamily\bfseries\footnotesize, text=bkRed] at (77.5,13) {2こ とる};
  \draw[bk arrow, bkTheme] (84,5) -- (91,5);
  % ③ 8
  \bktenframe{(94,0)}{5}{8}{bkRoute2}{0}{bkRoute2}
  \node[font=\sffamily\bfseries\small] at (17,-5) {13};
  \node[font=\sffamily\bfseries\small] at (67.5,-5) {10};
  \node[font=\sffamily\bfseries\small] at (106.5,-5) {8};
\end{tikzpicture}
\end{bkfigure}

\begin{enumerate}[label={\bkcircled[bkTheme]{\arabic*}}, leftmargin=2\zw, itemsep=0pt, topsep=0pt]
\item まず ばらの 3こを とる。 13 － 3 ＝ 10。
\item のこりの 2こを、 10から とる。 10 － 2 ＝ 8。
\end{enumerate}
\noindent どちらの やりかたでも 8に なりました。 じぶんが わかりやすい ほうを つかおう。

\noindent こたえが あって いるかは、 たし\R{算}{ざん}で たしかめられます。 のこった 8こに、 とった 5こを もどすと、 はじめの 13こに もどる はずです。 8 ＋ 5 ＝ 13。 あって いるね。

\bksubsection{5}{0が でて くる ひき\R{算}{ざん}}

\label{bk:1-5:zero}クッキーが 3まい。 3まい ぜんぶ たべると、 のこりは なんまい？

\noindent\begin{minipage}[c]{0.44\linewidth}\centering
\begin{tikzpicture}[x=1mm, y=1mm]
  \bkplate{(16,0)}{36}
  \foreach \x in {6,16,26} {\bkqegone{(\x,1.5)}{3.6}}
  \node[font=\sffamily\bfseries\small] at (16,-8) {のこり 0まい};
\end{tikzpicture}
\end{minipage}\hfill
\begin{minipage}[c]{0.54\linewidth}
ぜんぶ たべたので、 なにも のこって いません。 なにも ない ことは 0（1-3）。
\end{minipage}
\bkbigshiki{3 － 3 ＝ 0}{3 ひく 3 は 0}
\noindent おなじ かずを ひくと、 0に なります。 また、 1まいも たべなければ、 へらないので 3 － 0 ＝ 3。 0を ひいても、 かずは かわりません。

\begin{taisetsu}
「のこり」「ちがい」「かわった \R{量}{りょう}」は、 どれも ひき\R{算}{ざん}。\\
13 － 5 ＝ 8 は、 8 ＋ 5 ＝ 13 で たしかめよう。
\end{taisetsu}

%% （第11回：ここで ページを かえない）
%% ---------------- p.60〜61 例題 ----------------

\begin{reidai}[のこり]
あめが 8こ。 3こ たべると、 のこりは なんこ？

\noindent\begin{minipage}[c]{0.66\linewidth}
\kangaekata へるから、 ひき\R{算}{ざん}。

\shiki 8 － 3 ＝ 5 \kotaenext 5こ
\end{minipage}\hfill
\begin{minipage}[c]{0.32\linewidth}\centering
\begin{tikzpicture}[x=1mm, y=1mm]
  \foreach \x in {0,7.5,15,22.5,30} {\bkcandy{(\x,0)}{1.8}{bkRoute5}}
  \foreach \x in {7.5,15,22.5} {\bkqegone{(\x,-8)}{2.8}}
\end{tikzpicture}
\end{minipage}
\end{reidai}

\begin{jibun}
すずめが 7わ。 3わ とんで いくと、 のこりは なんわ？
\bkshikiran
\bkjans{7 － 3 ＝ 4\quad 4わ}
\end{jibun}

\begin{reidai}[ちがい]
あかい くるまが 7だい、 あおい くるまが 4だい。 どちらが なんだい おおい？
\begin{bkfigure}
\begin{tikzpicture}[x=1mm, y=1mm]
  \foreach \x in {0,12,24,36,48,60,72} {\bkcar{(\x,9)}{0.3}{bkAppleRed}}
  \foreach \x in {0,12,24,36} {\bkcar{(\x,0)}{0.3}{bkRoute2}}
  \foreach \x in {3.6,15.6,27.6,39.6} {\draw[bkGreen, line width=1.1pt] (\x,8.4) -- (\x,5.6);}
  \draw[bkOrange, line width=1pt, dash pattern=on 2pt off 1.5pt, rounded corners=2mm] (45,7.2) rectangle (82,15);
\end{tikzpicture}
\end{bkfigure}
\kangaekata ペアに すると、 あかが 3だい あまる。

\shiki 7 － 4 ＝ 3 \kotaenext あかが 3だい おおい
\end{reidai}

\begin{jibun}
いぬが 9ひき、 ねこが 5ひき。 どちらが なんびき おおい？
\bkshikiran[11\zw]
\bkjans{9 － 5 ＝ 4\quad いぬが 4ひき おおい}
\end{jibun}

%% （第11回：ここで ページを かえない）
\begin{reidai}[かわった \R{量}{りょう}]
ひまわりの たかさが、 ものさしの \R{目}{め}もり 5から 9に なりました。 なん\R{目}{め}もり のびた？

\kangaekata どれだけ かわったかは、 あと － まえ。

\shiki 9 － 5 ＝ 4 \kotaenext 4\R{目}{め}もり
\end{reidai}

\begin{jibun}
バスに 3\R{人}{にん} のって いました。 つぎに \R{見}{み}ると 8\R{人}{にん}。 なん\R{人}{にん} ふえた？
\bkshikiran
\bkjans{8 － 3 ＝ 5\quad \mbox{5\R{人}{にん}}}
\end{jibun}

\begin{reidai}[10の まとまりを くずす]
たまごが 12こ。 7こ つかうと、 のこりは なんこ？

\kangaekata 12は 10と 2。 2から 7は とれないので、 10から 7を とって 3。 3と 2で 5。

\shiki 12 － 7 ＝ 5 \kotaenext 5こ

\bkstep[bkGreen]{たしかめ}5 ＋ 7 ＝ 12
\end{reidai}

\begin{jibun}
いろがみが 14まい。 6まい つかうと、 のこりは なんまい？\enspace{\small\color{bkOrange!58!black}ヒント：10から 6を とると 4}
\bkshikiran
\bkjans{14 － 6 ＝ 8\quad 8まい}
\end{jibun}

%% ---------------- 第11回：0が でて くる ひき算 ----------------
\begin{reidai}[0が でて くる ひき\R{算}{ざん}]
（1）ふうせんが 4こ。 4こ とんで いくと、 のこりは なんこ？\par
（2）ふうせんが 4こ。 1こも とんで いかないと、 なんこ？

\kangaekata （1）ぜんぶ なくなるので 0。（2）へらないので 4の まま。

\shiki （1）4 － 4 ＝ 0\enspace（2）4 － 0 ＝ 4 \kotaenext （1）0こ\enspace（2）4こ
\end{reidai}

\begin{jibun}
にわに すずめが 6わ。 6わ とんで いくと、 のこりは なんわ？
\bkshikiran
\bkjans{6 － 6 ＝ 0\quad 0わ}
\end{jibun}

\begin{jibun}[どの ひき\R{算}{ざん}かな？]
しきと こたえを かこう。 「のこり」「ちがい」「かわった \R{量}{りょう}」の どれかも かんがえよう。\par\vspace{2pt}
（1）みかんが 9こ。 4こ たべると、 のこりは なんこ？
\bkshikiran
（2）えんぴつが 11\R{本}{ぽん}、 けしゴムが 6こ。 どちらが いくつ おおい？
\bkshikiran[11\zw]
（3）ひまわりの たかさが、 \R{目}{め}もり 6から 13に なった。 なん\R{目}{め}もり のびた？
\bkshikiran
\bkjans{（1）9 － 4 ＝ 5、 5こ（のこり） （2）11 － 6 ＝ 5、 えんぴつが 5つ おおい（ちがい） （3）13 － 6 ＝ 7、 7\R{目}{め}もり（かわった \R{量}{りょう}）}
\end{jibun}

\begin{column}{3 － 5 は できないの？}
おはじき 3こから、 5こは とれません。 でも、 0より \R{小}{ちい}さい かずを つかうと、 3 － 5 も けいさんできます。 さむい \R{日}{ひ}の おんどけいには、 0より \R{下}{した}の \R{目}{め}もりが あるよ。 この \R{本}{ほん}の 2\R{巻}{かん}で まなびます。
\end{column}

\clearpage
%% ---------------- p.62 まちがえやすいところ ----------------
\bksection{tsumazuki}{まちがえやすい ところ}

\begin{tsumazuki}{「ちがい」なのに たし\R{算}{ざん}？}
あかい はな 7\R{本}{ほん}と、 しろい はな 4\R{本}{ほん}の ちがいは？\\
\bkbatsu\enspace 7 ＋ 4 ＝ 11\hspace{3\zw}\bkmaru\enspace 7 － 4 ＝ 3\par
「ちがいは」「どちらが おおい」は、 ひき\R{算}{ざん}です。
\end{tsumazuki}

\begin{tsumazuki}{「ふえた」から たし\R{算}{ざん}？}
どんぐりが 3こから 8こに なりました。 なんこ ふえた？\\
\bkbatsu\enspace 3 ＋ 8 ＝ 11\hspace{3\zw}\bkmaru\enspace 8 － 3 ＝ 5\par
「どれだけ ふえた？」は、 あと － まえ。 ひき\R{算}{ざん}です。
\end{tsumazuki}

\begin{tsumazuki}{0を ひくと 0？}
\bkbatsu\enspace 7 － 0 ＝ 0\hspace{3\zw}\bkmaru\enspace 7 － 0 ＝ 7\par
0を ひくのは、 1つも とらない こと。 7の ままです。
\end{tsumazuki}

\begin{tsumazuki}{5 － 3 を して しまう}
\bkbatsu\enspace 13 － 5 ＝ 12\hspace{3\zw}\bkmaru\enspace 13 － 5 ＝ 8\par
ばらの 3から 5は とれない。 10の まとまりから とろう。
\end{tsumazuki}



\clearpage
%% ---------------- p.63 つかわれているところ ----------------
\begin{butsuri}
ひき\R{算}{ざん}は、 まちの なかで はたらいて います。
\begin{bkfigure}
\begin{tikzpicture}[x=1mm, y=1mm]
  % レジの おつり
  \fill[white, rounded corners=3mm] (0,0) rectangle (59,40);
  \fill[bkGray!70, rounded corners=1.5mm] (8,13) rectangle (36,22);
  \foreach \x in {13,20,27,34} {\fill[white] (\x-2.4,15.5) rectangle ++(4,3);}
  \fill[bkInk!85, rounded corners=1mm] (6,24) rectangle (38,36);
  \node[text=bkYellow, font=\sffamily\bfseries\footnotesize] at (22,32.5) {100 － 70};
  \node[text=bkYellow, font=\sffamily\bfseries\footnotesize] at (22,27.5) {おつり 30 \R{円}{えん}};
  \foreach \y in {14,22,30} {\bkcoin{(48,\y)}{10}}
  \node[font=\sffamily\bfseries\small] at (29.5,5.5) {レジの おつり};
  % スピードメーター
  \begin{scope}[shift={(63,0)}]
  \fill[white, rounded corners=3mm] (0,0) rectangle (59,40);
  \bkqemeter{(17,24)}{12}{50}
  \bkcar{(36,15)}{0.55}{bkRoute6}
  \draw[bk arrow, bkRoute6] (36,11) -- (54,11);
  \node[font=\sffamily\bfseries\small] at (29.5,5.5) {スピードメーター};
  \end{scope}
\end{tikzpicture}
\end{bkfigure}
\begin{itemize}[label={\color{bkPurple}●}, leftmargin=1.4\zw, itemsep=0pt, topsep=0pt]
\item レジは、 だした おかねから ねだんを ひいて、 おつりを だします。 100\R{円}{えん}は 10\R{円}{えん}\R{玉}{だま} 10まい（1-2）。 70\R{円}{えん}の ものなら、 10 － 7 ＝ 3で、 おつりは 10\R{円}{えん}\R{玉}{だま} 3まい。
\item スピードメーターは、 すこしの あいだに タイヤが どれだけ まわったかで、 はやさを だします。
\item おんどけいで、 あさより どれだけ あたたかく なったかが わかります。
\end{itemize}
\end{butsuri}

\bkfukidashi{delta}{「どれだけ かわったか」が わかると、 はやさや そだちかたも わかるんだ。 「\R{変化}{へん|か}の \R{路線}{ろ|せん}」へ しゅっぱつ！}

\begin{bkfigure}
\begin{tikzpicture}[x=1mm, y=1mm]
  \draw[bkRoute1, line width=3pt] (10,0) -- (128,0);
  \node[anchor=west, inner sep=0pt] at (0,12) {\bkrouteicon[1.4em]{1}\hspace{0.3em}{\sffamily\bfseries\small\color{bkRoute1!80!black}\R{変化}{へん|か}の \R{路線}{ろ|せん}}};
  \foreach \x/\c/\t/\s in {10/bkVol1/{1-5}/{ひき\R{算}{ざん}},39.5/bkVol1/{1-10}/{\R{時間}{じ|かん}},69/bkVol1/{1-14}/{1つ\R{分}{ぶん}},98.5/bkVol1/{1-18}/{おれ\R{線}{せん}グラフ},128/bkVol2/{2-3}/{{\ltjsetparameter{jacharrange={+200}}＋－×÷}}} {
    \fill[white, draw=\c, line width=1.6pt] (\x,0) circle (2.6);
    \node[anchor=base, font=\sffamily\bfseries\footnotesize, text=\c!70!black] at (\x,-8.4) {\t};
    \node[anchor=base, font=\sffamily\footnotesize, text=bkInk] at (\x,-15.8) {\s};}
  \node[fill=bkTheme, text=white, rounded corners=1mm, font=\sffamily\bfseries\fontsize{8}{9}\selectfont, inner sep=1.5pt] at (10,5.5) {いま ここ};
\end{tikzpicture}
\end{bkfigure}

\clearpage
%% ---------------- p.64〜65 もんだい ----------------
\bkendproblems

\bkproblemgroup{きほん}{1}

\bkproblem けいさんを しましょう。
\par\vspace{2pt}
\noindent\hspace*{1\zw}\begin{tblr}{colspec={*{2}{Q[l,m,wd=62mm]}}, colsep=0pt, rowsep=3pt}
（1）6 － 2 ＝ \bkkakikomi & （2）9 － 4 ＝ \bkkakikomi \\
（3）10 － 4 ＝ \bkkakikomi & （4）7 － 5 ＝ \bkkakikomi \\
（5）12 － 9 ＝ \bkkakikomi & （6）15 － 7 ＝ \bkkakikomi \\
（7）16 － 8 ＝ \bkkakikomi & （8）11 － 4 ＝ \bkkakikomi \\
（9）6 － 6 ＝ \bkkakikomi & （10）8 － 0 ＝ \bkkakikomi \\
\end{tblr}
\bkpans{（1）4\quad（2）5\quad（3）6\quad（4）2\quad（5）3\quad（6）8\quad（7）8\quad（8）7\quad（9）0\quad（10）8}

\bkproblem ふうせんが 7こ。 えを みて、 ひき\R{算}{ざん}の しきを かこう。
\begin{bkfigure}
\begin{tikzpicture}[x=1mm, y=1mm]
  \fill[bkTheme!6, rounded corners=3mm] (0,-12) rectangle (120,16);
  \foreach \x in {10,20,30,40} {\bkballoon{(\x,2)}{3}{bkAppleRed}}
  \foreach \x/\y in {76/6,88/10,100/7} {\bkballoon{(\x,\y)}{3}{bkAppleRed}}
  \draw[bk arrow, bkGray] (54,0) -- (66,8);
  \node[font=\sffamily\bfseries\small, text=bkInk] at (88,-6) {3こ とんで いった};
\end{tikzpicture}
\end{bkfigure}
\bkshikiran
\bkpans{7 － 3 ＝ 4\quad 4こ}

\bkproblem きんぎょ 8ひき、 めだか 5ひき。 どちらが なんびき おおい？
\bkshikiran[11\zw]
\bkpans{8 － 5 ＝ 3\quad きんぎょが 3びき おおい}

\bkproblem シールが 4まいから 9まいに なった。 なんまい ふえた？
\bkshikiran
\bkpans{9 － 4 ＝ 5\quad 5まい}

\bkproblemgroup{チャレンジ}{1}

\bkproblem わなげで、 ソラは 13てん、 ミナは 8てん。 どちらが なんてん おおい？
\bkshikiran[11\zw]
\bkpans{13 － 8 ＝ 5\quad ソラが 5てん おおい}

\bkproblem あさがおの たかさが、 \R{目}{め}もり 9から 14に なりました。 なん\R{目}{め}もり のびた？ こたえを たし\R{算}{ざん}で たしかめよう。
\bkshikiran
\par\vspace{4pt}\noindent\tikz[baseline=(s.base)]\node[fill=bkGreen, text=white, rounded corners=1.2mm, font=\sffamily\bfseries\small, inner xsep=5pt, inner ysep=2.5pt] (s) {たしかめ};\enspace\bkkakikomi[60mm]
\bkpans{14 － 9 ＝ 5、 5\R{目}{め}もり （たしかめ 5 ＋ 9 ＝ 14）}
\par\vspace{5mm}
\begin{furikaeri}
\begin{checklist}
\item 「のこり」「ちがい」を しきに かける。\hfill{\footnotesize\sffamily\color{bkTheme}→\pageref{bk:1-5:nokori}ページ}
\item どれだけ かわったかを もとめられる。\hfill{\footnotesize\sffamily\color{bkTheme}→\pageref{bk:1-5:kawatta}ページ}
\item 10の まとまりを くずして ひける。\hfill{\footnotesize\sffamily\color{bkTheme}→\pageref{bk:1-5:kuzusu}ページ}
\end{checklist}
チェックが つかなかったら、 →の ページに もどろう。
\end{furikaeri}

\noindent\begin{minipage}[c]{\dimexpr\linewidth-60mm\relax}
\bkfukidashi{sora}{あしたも はかって、 きょうと くらべよう！}
\end{minipage}\hfill
\begin{minipage}[c]{57.5mm}\scalebox{0.85}{\bkkigocard{－}{ひく}{のこりや ちがいを もとめる しるし}{minus}{1-5}}\end{minipage}
\clearpage
